% 공통 프리앰블: 서강대 PHY1001 다이어그램
\documentclass[tikz,border=20pt]{standalone}
\usepackage{kotex}
\usepackage{amsmath}
\usepackage{tikz}
\usetikzlibrary{arrows.meta, decorations.pathmorphing, decorations.markings, patterns, calc, angles, quotes, positioning}

% 서강대 색상 팔레트
\definecolor{sgred}{HTML}{8B0029}
\definecolor{sgblue}{HTML}{2563EB}
\definecolor{sggray}{HTML}{6B7280}
\definecolor{sglight}{HTML}{F3F4F6}

% 공통 스타일
\tikzset{
  vec/.style={-{Stealth[length=8pt,width=5pt]}, line width=2.5pt},
  vecred/.style={vec, sgred},
  vecblue/.style={vec, sgblue},
  auxline/.style={sggray, line width=1pt, dashed, dash pattern=on 4pt off 3pt},
  block/.style={draw=sggray, line width=2pt, fill=sglight, minimum width=1.2cm, minimum height=1.2cm},
  ground/.style={pattern=north east lines, pattern color=sggray!40},
  labelstyle/.style={font=\sffamily\fontsize{14}{16}\selectfont},
  mathlabel/.style={font=\sffamily\itshape\fontsize{14}{16}\selectfont},
  titlestyle/.style={font=\sffamily\bfseries\fontsize{16}{18}\selectfont, sgred},
  spring/.style={decorate, decoration={zigzag, segment length=6pt, amplitude=5pt}, line width=2pt, sggray},
}

\begin{document}
\begin{tikzpicture}

  % === (a) 평형 상태 ===
  \begin{scope}[shift={(0, 8)}]
    \node[labelstyle, sggray, font=\sffamily\bfseries] at (-1.2, 0.6) {(a)};

    % Wall
    \fill[ground] (-0.5, -0.4) rectangle (0, 1.2);
    \draw[sggray, line width=2pt] (0, -0.4) -- (0, 1.2);

    % Spring (relaxed)
    \draw[spring] (0, 0.4) -- (3.5, 0.4);

    % Block at x=0
    \node[block, minimum width=1.2cm, minimum height=1.0cm] (b1) at (4.1, 0.4) {};
    \node[labelstyle, sggray, font=\sffamily\fontsize{12}{14}\selectfont] at (b1.center) {$m$};

    % x-axis
    \draw[sggray, line width=1pt, -{Stealth[length=6pt]}] (-0.3, -0.7) -- (8, -0.7);
    \node[mathlabel, sggray] at (8.3, -0.7) {$x$};

    % Origin mark
    \draw[sggray, line width=1pt] (4.1, -0.5) -- (4.1, -0.9);
    \node[labelstyle, sggray] at (4.1, -1.2) {$0$};

    % Label
    \node[labelstyle, sgred, font=\sffamily\fontsize{13}{15}\selectfont] at (6.5, 0.4) {평형 상태, $F_s = 0$};
  \end{scope}

  % === (b) 늘어난 상태 (x > 0) ===
  \begin{scope}[shift={(0, 4)}]
    \node[labelstyle, sggray, font=\sffamily\bfseries] at (-1.2, 0.6) {(b)};

    % Wall
    \fill[ground] (-0.5, -0.4) rectangle (0, 1.2);
    \draw[sggray, line width=2pt] (0, -0.4) -- (0, 1.2);

    % Spring (stretched)
    \draw[spring, decoration={segment length=9pt}] (0, 0.4) -- (5, 0.4);

    % Block at x > 0
    \node[block, minimum width=1.2cm, minimum height=1.0cm] (b2) at (5.6, 0.4) {};
    \node[labelstyle, sggray, font=\sffamily\fontsize{12}{14}\selectfont] at (b2.center) {$m$};

    % Force arrow (left, restoring)
    \draw[vecred] (4.95, 0.4) -- (3.5, 0.4);
    \node[labelstyle, sgred, above=3pt] at (4.2, 0.5) {$\vec{F}_s$};

    % Displacement arrow (blue, from 0 to block)
    \draw[vecblue] (4.1, -0.3) -- (5.6, -0.3);
    \node[labelstyle, sgblue, below=2pt] at (4.85, -0.4) {$\vec{d}$};

    % x-axis
    \draw[sggray, line width=1pt, -{Stealth[length=6pt]}] (-0.3, -0.7) -- (8, -0.7);
    \node[mathlabel, sggray] at (8.3, -0.7) {$x$};

    % Equilibrium dashed line
    \draw[auxline] (4.1, -0.5) -- (4.1, 1.2);
    \node[labelstyle, sggray] at (4.1, -1.2) {$0$};

    % x position mark
    \draw[sggray, line width=1pt] (5.6, -0.5) -- (5.6, -0.9);
    \node[labelstyle, sgblue] at (5.6, -1.2) {$x > 0$};

    % Label
    \node[labelstyle, sgred, font=\sffamily\fontsize{13}{15}\selectfont] at (7, 1.1) {$F_s = -kx$ (복원력)};
    \node[labelstyle, sggray, font=\sffamily\fontsize{11}{13}\selectfont] at (7, 0.5) {← 평형 쪽으로};
  \end{scope}

  % === (c) 압축된 상태 (x < 0) ===
  \begin{scope}[shift={(0, 0)}]
    \node[labelstyle, sggray, font=\sffamily\bfseries] at (-1.2, 0.6) {(c)};

    % Wall
    \fill[ground] (-0.5, -0.4) rectangle (0, 1.2);
    \draw[sggray, line width=2pt] (0, -0.4) -- (0, 1.2);

    % Spring (compressed)
    \draw[spring, decoration={segment length=3.5pt}] (0, 0.4) -- (2.0, 0.4);

    % Block at x < 0
    \node[block, minimum width=1.2cm, minimum height=1.0cm] (b3) at (2.6, 0.4) {};
    \node[labelstyle, sggray, font=\sffamily\fontsize{12}{14}\selectfont] at (b3.center) {$m$};

    % Force arrow (right, restoring)
    \draw[vecred] (3.25, 0.4) -- (4.7, 0.4);
    \node[labelstyle, sgred, above=3pt] at (4.0, 0.5) {$\vec{F}_s$};

    % Displacement arrow (blue, from 0 to block, leftward)
    \draw[vecblue] (4.1, -0.3) -- (2.6, -0.3);
    \node[labelstyle, sgblue, below=2pt] at (3.35, -0.4) {$\vec{d}$};

    % x-axis
    \draw[sggray, line width=1pt, -{Stealth[length=6pt]}] (-0.3, -0.7) -- (8, -0.7);
    \node[mathlabel, sggray] at (8.3, -0.7) {$x$};

    % Equilibrium dashed line
    \draw[auxline] (4.1, -0.5) -- (4.1, 1.2);
    \node[labelstyle, sggray] at (4.1, -1.2) {$0$};

    % x position mark
    \draw[sggray, line width=1pt] (2.6, -0.5) -- (2.6, -0.9);
    \node[labelstyle, sgblue] at (2.6, -1.2) {$x < 0$};

    % Label
    \node[labelstyle, sgred, font=\sffamily\fontsize{13}{15}\selectfont] at (7, 1.1) {$F_s = -kx$ (복원력)};
    \node[labelstyle, sggray, font=\sffamily\fontsize{11}{13}\selectfont] at (7, 0.5) {→ 평형 쪽으로};
  \end{scope}

\end{tikzpicture}
\end{document}
